%% Aplikovaná dátová analýza -- Sylabus, zimný semester 2026 (slovenská verzia)
%% Kompilácia: pdflatex syllabus_ADA_2026_sk.tex
\documentclass[11pt]{article}

\usepackage[a4paper,top=2.2cm,bottom=2.2cm,left=2.2cm,right=2.2cm]{geometry}
\usepackage[T1]{fontenc}
\usepackage[utf8]{inputenc}
\usepackage[slovak]{babel}
\usepackage{lmodern}
\usepackage{microtype}
\usepackage{array}
\usepackage{longtable}
\usepackage{booktabs}
\usepackage{enumitem}
\usepackage{parskip}
\usepackage{titlesec}
\usepackage{xcolor}
\usepackage[hidelinks]{hyperref}

\definecolor{oxfordblue}{HTML}{002147}

\titleformat{\section}{\large\bfseries\color{oxfordblue}}{}{0pt}{\MakeUppercase}
\titlespacing*{\section}{0pt}{12pt}{4pt}

\newcommand{\hdivider}{\noindent\rule{\textwidth}{0.8pt}}
\newcommand{\readings}[1]{\par\vspace{3pt}{\footnotesize\itshape #1\par}}

\setlist[itemize]{leftmargin=1.4em,topsep=2pt,itemsep=1pt}
\setlength{\LTpre}{4pt}
\setlength{\LTpost}{4pt}

\newcolumntype{W}{>{\raggedright\arraybackslash}p{0.8cm}}
\newcolumntype{T}{>{\raggedright\arraybackslash}p{10.6cm}}
\newcolumntype{D}{>{\raggedright\arraybackslash}p{4.1cm}}

\pagestyle{plain}

\begin{document}

%% ---------------------------------------------------------------- HLAVIČKA
\begin{center}
{\LARGE\bfseries\color{oxfordblue} APLIKOVANÁ DÁTOVÁ ANALÝZA}\\[2pt]
{\large Národohospodárska fakulta}\\[1pt]
{\large Ekonomická univerzita v Bratislave}\\[1pt]
{\large Katedra hospodárskej politiky}
\end{center}

\hdivider
\smallskip

\noindent
\begin{tabular}{@{}>{\bfseries}p{3.4cm}p{12.6cm}@{}}
Vyučujúci: & Tomáš Oleš, Martin Vaško, Lucia Širochmanová, Ella Sargsyan \\
E-mail: & tomas.oles@euba.sk; martin.vasko@euba.sk; lucia.sirochmanova@euba.sk;\newline esargsyan@cerge-ei-foundation.org \\
Konzultačné hodiny: & streda 13:00 -- 15:00 \\
Kancelária: & Stará budova, 4B40 \\
Kredity: & 6 \\
Semester: & zimný semester 2026 \\
\end{tabular}

\smallskip
\hdivider

%% ---------------------------------------------------------------- OPIS
\section{Charakteristika predmetu}

Prvým cieľom predmetu je poskytnúť študentom vedomosti a zručnosti v moderných metódach aplikovanej dátovej
analýzy a štatistického učenia vrátane práce v prostredí R, aby dokázali realizovať empirický ekonomický
výskum a navrhovať výskumné prístupy a metódy na riešenie ekonomických problémov.

Druhým cieľom predmetu je oboznámiť študentov so základnými ekonometrickými pojmami a metódami, ktoré sú
nevyhnutné na čítanie a porozumenie súčasnému empirickému výskumu v ekonómii, ako aj na plánovanie a
realizáciu samostatných výskumných projektov. Predmet uvedie myšlienku regresnej analýzy a pokryje niektoré
z najnovších ekonometrických techník, ktoré tvoria jadro modernej ekonometrickej praxe zameranej na
identifikáciu kauzálnych vzťahov.

Predmet je rozdelený na dve časti. \textbf{Časť I -- Aplikovaná dátová analýza v R} (týždne 1--7) pokrýva
dátovo-analytický nástrojový aparát: jazyk R, vizualizáciu údajov, spracovanie údajov, analýzu textu a úvod
do štatistického učenia prostredníctvom klasifikácie a zhlukovania. \textbf{Časť II -- Ekonometria v kocke}
(týždne 8--13) je samostatný modul vedený Ellou Sargsyan, Ph.D.\ (CERGE-EI Foundation) a venuje sa regresnej
analýze a moderným technikám kauzálnej inferencie.

%% ---------------------------------------------------------------- VÝSTUPY
\section{Výstupy vzdelávania}

Po absolvovaní predmetu študenti získajú poznatky o moderných metódach dátového výskumu a vizualizácie,
klasifikácie, zhlukovania, lineárnej regresie a všeobecnej dátovej analýzy. Získajú zručnosti v práci s
údajmi, ktoré budú vedieť uplatniť vo vlastnom empirickom výskume. Rovnako nadobudnú pokročilé zručnosti v
používaní moderného softvéru (R) v empirickom ekonomickom výskume a budú schopní písať vlastné funkcie.

Študenti budú poznať základné ekonometrické pojmy, základné metódy odhadu a metódy testovania štatistických
hypotéz. Budú schopní aplikovať štandardné postupy konštrukcie ekonometrických modelov, spracovať
štatistické informácie, dospieť k štatisticky podloženým záverom a zmysluplne interpretovať výsledky
odhadnutých ekonometrických modelov. Okrem toho študenti získajú praktické zručnosti pri spracovaní
reálnych údajov s využitím ekonometrických balíkov na tvorbu a odhad regresných modelov v R.

%% ---------------------------------------------------------------- PREREKVIZITY
\section{Prerekvizity}

Úvod do štatistiky.

%% ---------------------------------------------------------------- LITERATÚRA
\section{Literatúra a študijné materiály}

\begin{itemize}
  \item Wickham, H., Çetinkaya-Rundel, M., Grolemund, G. (2023). \textit{R for Data Science}. O'Reilly
        Media, Inc. Voľne dostupné na \url{https://r4ds.hadley.nz/}.
  \item Gareth, J., Daniela, W., Trevor, H., Robert, T. (2013). \textit{An Introduction to Statistical
        Learning: With Applications in R}. Springer. Voľne dostupné na \url{https://www.statlearning.com/}.
  \item Silge, J., Robinson, D. (2017). \textit{Text Mining with R: A Tidy Approach}. O'Reilly Media, Inc.
        Voľne dostupné na \url{https://www.tidytextmining.com/}.
  \item Wooldridge, Jeffrey M. \textit{Introductory Econometrics: A Modern Approach}. 5.~vyd., US, 2012.
  \item Angrist, J. D., Pischke, J.-S. (2009). \textit{Mostly Harmless Econometrics: An Empiricist's
        Companion}. Princeton University Press.
  \item Tutoriál k základom R: \url{https://cran.r-project.org/doc/contrib/Paradis-rdebuts_en.pdf}.
\end{itemize}

Všetky študijné materiály (príslušné kapitoly učebníc, odborné články a prezentácie z prednášok) budú
zverejnené na stránke predmetu v systéme Moodle.

%% ---------------------------------------------------------------- HODNOTENIE
\section{Hodnotenie predmetu}

Výsledná známka sa skladá z troch zložiek.

\begin{itemize}
  \item \textbf{Časť I --- kvízy a domáce zadania: 30\,\%.} Krátke kvízy a zadania zadávané počas prednášok
        Časti I aj po nich, spolu s dvoma domácimi zadaniami v prostredí R.
  \item \textbf{Časť II --- Ekonometria v kocke: 35\,\%.} Hodnotí sa v rámci tohto modulu.
  \item \textbf{Záverečný dátový projekt: 35\,\%.} Dátový projekt prezentovaný na konci semestra. Vychádza
        z Časti I aj Časti II a je záverečnou prácou predmetu.
\end{itemize}

\noindent\textbf{Otázky a diskusia na hodine} sa bodovo nehodnotia, aktívna účasť však môže byť zohľadnená
podľa uváženia vyučujúceho.

\smallskip
\noindent\textbf{Výpočet výslednej známky:} záverečná písomná skúška sa nekoná. Študenti nemusia dosiahnuť
stanovenú hranicu v každej jednotlivej zložke; výsledná známka vychádza z kumulatívneho skóre.

%% ---------------------------------------------------------------- TERMÍNY
\section{Dôležité termíny}

\begin{itemize}
  \item \textbf{Časť I (Aplikovaná dátová analýza v R):} 22. september -- 4. november 2026
  \item \textbf{Časť II (Ekonometria v kocke):} 26. október -- 4. december 2026
  \item \textbf{Prezentácie záverečných projektov:} 7. -- 11. december 2026
  \item \textbf{Opravný termín:} 14. -- 18. december 2026
\end{itemize}

\noindent\textbf{Časy výučby.}

\smallskip
\noindent\begin{tabular}{@{}>{\bfseries}p{2.4cm}p{6.4cm}p{6.4cm}@{}}
\toprule
 & \textbf{Časť I} --- Aplikovaná dátová analýza v R & \textbf{Časť II} --- Ekonometria v kocke \\
\midrule
Kód predmetu & NND21210/21 (slovenská skupina) & modul CERGE-EI \\
Deň & utorok, od 22. 9. 2026 & Prednáška: po + št, od 26. 10. 2026\newline Cvičenie: utorok \\
Prednáška & 13:30--15:00, miestnosť 4B26 & 11:00--12:30 SEČ, \textbf{online} \\
Cvičenie & 15:00--16:30, miestnosť 4B26 & 13:30--15:00, miestnosť 4B26 (prezenčne) \\
\bottomrule
\end{tabular}

\smallskip
\noindent Tento sylabus uvádza časy pre \textbf{slovenskú skupinu}. Anglická skupina
(NND21254/21) má výučbu v stredu --- pozri anglickú verziu sylabu.

\smallskip
\noindent V Časti II sú prednášky vedené \textbf{online} (Ella Sargsyan), cvičenie však
zostáva \textbf{prezenčné} v miestnosti 4B26 v obvyklom utorkovom čase.

\smallskip
\noindent\footnotesize Vyučujúci si vyhradzuje právo zmeniť obsah študijného materiálu, ak to bude
považovať za potrebné z dôvodu odloženia výučby (nepriaznivé počasie, choroba vyučujúceho a pod.) alebo z
dôvodu tempa výučby.\normalsize

%% ---------------------------------------------------------------- INTEGRITA
\section{Akademická integrita}

Akademická integrita je nevyhnutným predpokladom vzdelávania a vedeckého bádania na univerzite a zárukou
toho, že diplom Ekonomickej univerzity v Bratislave je silným signálom individuálneho akademického výkonu
každého študenta. Univerzita preto pristupuje k prípadom podvádzania a plagiátorstva veľmi vážne. Etický
kódex Ekonomickej univerzity v Bratislave (\url{https://euba.sk/univerzita/eticky-kodex}) vymedzuje
konanie, ktoré predstavuje akademickú nečestnosť, a postupy riešenia akademických previnení. Medzi možné
previnenia patria okrem iného:

\smallskip
\noindent\textit{V prácach a zadaniach:} použitie myšlienok alebo slov inej osoby bez náležitého uvedenia
zdroja; odovzdanie vlastnej práce vo viac než jednom predmete bez súhlasu vyučujúceho; vymýšľanie zdrojov
alebo faktov; získanie alebo poskytnutie nepovolenej pomoci pri akomkoľvek zadaní.

\noindent\textit{Pri testoch a skúškach:} používanie alebo držba nepovolených pomôcok; nazeranie do
odpovedí inej osoby počas skúšky alebo testu; skresľovanie vlastnej identity či vydávanie sa za inú osobu.

\noindent\textit{V akademickej práci:} falšovanie dokumentov alebo známok; falšovanie alebo pozmeňovanie
akejkoľvek dokumentácie vyžadovanej univerzitou vrátane (nie však výlučne) lekárskych potvrdení.

\smallskip
Všetky podozrenia z akademickej nečestnosti budú prešetrené podľa postupov uvedených v Etickom kódexe. Ak
majú študenti otázky alebo pochybnosti o tom, čo predstavuje primerané akademické správanie alebo primerané
výskumné a citačné metódy, očakáva sa, že si vyhľadajú ďalšie informácie o akademickej integrite u svojich
vyučujúcich alebo z iných inštitucionálnych zdrojov.

\clearpage

%% ---------------------------------------------------------------- HARMONOGRAM
\section{Harmonogram výučby}

\begin{longtable}{@{}WTD@{}}
\toprule
\textbf{Tž.} & \textbf{Obsah témy} & \textbf{Termín a vyučujúci} \\
\midrule
\endfirsthead
\toprule
\textbf{Tž.} & \textbf{Obsah témy} & \textbf{Termín a vyučujúci} \\
\midrule
\endhead
\bottomrule
\endfoot

\multicolumn{3}{@{}l}{\textbf{\color{oxfordblue} ČASŤ I --- APLIKOVANÁ DÁTOVÁ ANALÝZA V R}}\\
\midrule

1 & \textbf{Základy R.} Základná syntax R, dátové typy, aritmetika vektorov, indexovanie, triedenie,
triedenie pomocou \texttt{dplyr} a vykresľovanie grafov v základných balíkoch. GitHub.
\readings{Literatúra: Gareth et al. (2013), kap.~1. Wickham et al. (2023), kap.~1, 2, 4--6.}
& utorok 22. 9. 2026\newline Prednáška 13:30--15:00, 4B26\newline Cvičenie 15:00--16:30, 4B26\newline \textit{L. Širochmanová} \\
\midrule

2 & \textbf{Vizualizácia -- časť 1.} Princípy vizualizácie údajov, tvorba vlastných grafov pomocou
\texttt{ggplot2} a rozbor výhod i úskalí najčastejšie používaných typov grafov.
\readings{Literatúra: Wickham et al. (2023), kap.~3, 7.}
& utorok 29. 9. 2026\newline Prednáška 13:30--15:00, 4B26\newline Cvičenie 15:00--16:30, 4B26\newline \textit{T. Oleš} \\
\midrule

3 & \textbf{Vizualizácia -- časť 2.} Princípy priestorovej vizualizácie, tvorba vlastných máp pomocou
\texttt{tmap}. Prehľad priestorových údajov (body, línie a polygóny) a stručná ukážka práce s rastrom.
\readings{Literatúra: Wickham et al. (2023), kap.~3, 7.}
& utorok 6. 10. 2026\newline Prednáška 13:30--15:00, 4B26\newline Cvičenie 15:00--16:30, 4B26\newline \textit{T. Oleš} \\
\midrule

4 & \textbf{Sumarizácia, upratovanie a spájanie údajov.} Import údajov z rôznych formátov, získavanie
údajov z webu, usporiadané (tidy) údaje s \texttt{tidyverse}, spracovanie údajov pomocou \texttt{dplyr},
práca s dátumovými a časovými formátmi.
\readings{Literatúra: Wickham et al. (2023), kap.~9--16.}
& utorok 13. 10. 2026\newline Prednáška 13:30--15:00, 4B26\newline Cvičenie 15:00--16:30, 4B26\newline \textit{T. Oleš} \\
\midrule

5 & \textbf{Analýza textu.} Usporiadaný textový formát a tokenizácia. Stop slová a frekvencie slov.
Analýza sentimentu na usporiadaných údajoch (lexikóny Bing, AFINN, NRC). Frekvencia termínu, inverzná
dokumentová frekvencia a \texttt{tf-idf}.
\readings{Literatúra: Silge a Robinson (2017), kap.~1--3.}
& utorok 20. 10. 2026\newline Prednáška 13:30--15:00, 4B26\newline Cvičenie 15:00--16:30, 4B26\newline \textit{T. Oleš} \\
\midrule

6--7\textsuperscript{**} & \textbf{Klasifikácia a zhluková analýza.}\textsuperscript{**} Úvod do
klasifikácie kvalitatívnych premenných. Klasifikátory: logistická regresia, naivný Bayesov klasifikátor,
metóda k-najbližších susedov, zovšeobecnené lineárne modely. Pretrénovanie, matica zámen, špecificita a
senzitivita. Učenie bez učiteľa a pojem podobnosti: metóda k-priemerov, hierarchické zhlukovanie a
dendrogramy, voľba počtu zhlukov, analýza hlavných komponentov a jej využitie v zhlukovej analýze.
\readings{Literatúra: Gareth et al. (2013), kap.~4 a kap.~12.}
& utorok 27. 10. a 3. 11. 2026\newline Prednáška 13:30--15:00, 4B26\newline Cvičenie 15:00--16:30, 4B26\newline \textit{T. Oleš, M. Vaško} \\
\midrule

\multicolumn{3}{@{}l}{\textbf{\color{oxfordblue} ČASŤ II --- EKONOMETRIA V KOCKE}}\\
\midrule

8 & \textbf{Podstata ekonometrie a ekonomických údajov.} Čo je ekonometria a prečo ju potrebujeme? Kroky
empirickej ekonomickej analýzy. Štruktúra ekonomických údajov. Kauzalita a pojem „ceteris paribus“ v
ekonometrických údajoch. Definícia jednoduchého regresného modelu.
\readings{Literatúra: Wooldridge (2012), kap.~1; kap.~2 (2.1).}
& \footnotesize\textbf{Prednáška (online)}\normalsize\newline pondelok 26. 10. 2026, 11:00--12:30\newline štvrtok 29. 10. 2026, 11:00--12:30\newline \footnotesize\textbf{Cvičenie (prezenčne)}\normalsize\newline utorok 27. 10. 2026, 13:30--15:00, 4B26\newline \textit{E. Sargsyan} \\
\midrule

9 & \textbf{Regresná analýza prierezových údajov: časť 1.} Metóda najmenších štvorcov. Vlastnosti metódy
najmenších štvorcov. Predpoklady jednoduchého regresného modelu.
\readings{Literatúra: Wooldridge (2012), kap.~2 (2.2--2.5).}
& \footnotesize\textbf{Prednáška (online)}\normalsize\newline pondelok 2. 11. 2026, 11:00--12:30\newline štvrtok 5. 11. 2026, 11:00--12:30\newline \footnotesize\textbf{Cvičenie (prezenčne)}\normalsize\newline utorok 3. 11. 2026, 13:30--15:00, 4B26\newline \textit{E. Sargsyan} \\
\midrule

10 & \textbf{Regresná analýza prierezových údajov: časť 2.} Viacnásobná regresná analýza. Testovanie
hypotéz. Prezentovanie výsledkov regresie.
\readings{Literatúra: Wooldridge (2012), kap.~3; kap.~4 (4.2--4.6).}
& \footnotesize\textbf{Prednáška (online)}\normalsize\newline pondelok 9. 11. 2026, 11:00--12:30\newline štvrtok 12. 11. 2026, 11:00--12:30\newline \footnotesize\textbf{Cvičenie (prezenčne)}\normalsize\newline utorok 10. 11. 2026, 13:30--15:00, 4B26\newline \textit{E. Sargsyan} \\
\midrule

11 & \textbf{Regresná analýza: ďalšie témy a problémy.} Merné jednotky a funkčné tvary. Regresná analýza s
kvalitatívnymi informáciami. Heteroskedasticita a sériová korelácia. Endogenita: vynechané premenné,
selekcia a spätná kauzalita.
\readings{Literatúra: Wooldridge (2012), kap.~2 (2.4); kap.~3 (3.3); kap.~7 (7.1--7.4); kap.~8.}
& \footnotesize\textbf{Prednáška (online)}\normalsize\newline pondelok 16. 11. 2026, 11:00--12:30\newline štvrtok 19. 11. 2026, 11:00--12:30\newline \footnotesize\textbf{Cvičenie (prezenčne)}\normalsize\newline utorok 17. 11. 2026, 13:30--15:00, 4B26\newline \textit{E. Sargsyan} \\
\midrule

12 & \textbf{Riešenie endogenity.} Inštrumentálne premenné. Dizajn regresnej diskontinuity.
\readings{Literatúra: Wooldridge (2012), kap.~15 (15.1--15.3); Angrist a Pischke (2009), kap.~6
(s.~189--195).\newline Články na diskusiu: Acemoglu, D., \& Johnson, S. (2005). „Unbundling institutions.“
\textit{Journal of Political Economy}, 113(5), 949--995. Carpenter, C., \& Dobkin, C. (2009). „The effect
of alcohol consumption on mortality: regression discontinuity evidence from the minimum drinking age.“
\textit{American Economic Journal: Applied Economics}, 1(1), 164--182.}
& \footnotesize\textbf{Prednáška (online)}\normalsize\newline pondelok 23. 11. 2026, 11:00--12:30\newline štvrtok 26. 11. 2026, 11:00--12:30\newline \footnotesize\textbf{Cvičenie (prezenčne)}\normalsize\newline utorok 24. 11. 2026, 13:30--15:00, 4B26\newline \textit{E. Sargsyan} \\
\midrule

13 & \textbf{Metódy panelových údajov.} Zavedenie časového rozmeru. Model fixných efektov. Metóda rozdielov
v rozdieloch.
\readings{Literatúra: Wooldridge (2012), kap.~14 (14.1); Angrist a Pischke (2009), kap.~5 (s.~169--182).
\newline Článok na diskusiu: Card, D., \& Krueger, A. B. (1994). „Minimum Wages and Employment: A Case
Study of the Fast-Food Industry in New Jersey and Pennsylvania.“ \textit{The American Economic Review},
84(4), 772--793.}
& \footnotesize\textbf{Prednáška (online)}\normalsize\newline pondelok 30. 11. 2026, 11:00--12:30\newline štvrtok 3. 12. 2026, 11:00--12:30\newline \footnotesize\textbf{Cvičenie (prezenčne)}\normalsize\newline utorok 1. 12. 2026, 13:30--15:00, 4B26\newline \textit{E. Sargsyan} \\

\end{longtable}

\smallskip
\noindent\footnotesize\textsuperscript{**} Preberá sa podľa úrovne a tempa skupiny. Tieto dve stretnutia
prebiehajú súbežne s úvodnými týždňami Časti II.\normalsize

\vfill
\begin{center}
\itshape „Sme tým, čo opakovane robíme. Dokonalosť teda nie je čin, ale zvyk.“ --- Aristoteles
\end{center}

\end{document}
